\section*{Chapter \ref{ch:b1:radiometric-methods} --- Radiometric Dating Methods}
% values: calc/ch25.py

\begin{solution}{exo:b1:radiometric-methods:1}
(a)~$D^*/P = 2^n - 1 = 3$ gives $n = 2$: two half-lives. (b)~Rubidium-87:
$e^{\lambda t} - 1 = e^{\num{1.3972e-11} \times \num{3.10e8}} - 1 = 0.004341$;
uranium-235: $e^{\num{9.8485e-10} \times \num{3.10e8}} - 1 = 0.3570$. The
uranium ratio is about 80 times larger: a given absolute error of measurement is
a far smaller part of it. (c)~$2^{-n} < 0.01$ when $n > \log_2 100 = 6.64$: after
about seven half-lives.
\end{solution}

\begin{solution}{exo:b1:radiometric-methods:2}
(a)~Radiocarbon on the charcoal. (b)~Uranium--lead on the zircon, or
argon--argon (potassium--argon) on the sanidine, a potassium-rich \omterm{def:b1:rock-forming-minerals:mineral}{mineral}.
(c)~Uranium--lead on zircon, which survives a long history; rubidium--strontium
on whole rocks is possible; potassium--argon would risk ages lowered by argon
lost in later heating. (d)~Radiocarbon on the carbonate of the shell, near the
limit of the method (about \num{70000} years). (e)~None directly: its age is
bracketed by dated volcanic rocks above and below, or found by correlating its
fossils with a dated section (\cref{ch:b1:stratigraphic-scale}).
\end{solution}

\begin{solution}{exo:b1:radiometric-methods:3}
(a)~$t = -8033 \ln 0.25 = 8033 \times 1.3863 = \num{11136}$ years BP.
(b)~The mean life is $5700/\ln 2 = \num{8223}$ years: $t = 8223 \times 1.3863 =
\num{11400}$ years. (c)~Both assume that the proportion of carbon-14 in living
carbon was the same then as in the standard; it was not. The calendar age is read
from a \omterm{def:b1:radiometric-methods:radiocarbon}{calibration curve} such as IntCal20, as a range of years cal BP.
\end{solution}

\begin{solution}{exo:b1:radiometric-methods:4}
(a)~Each clock starts when its \omterm{def:b1:rock-forming-minerals:mineral}{mineral} cools below its \omterm{def:b1:radiometric-methods:closure}{closure temperature}:
zircon (above \qty{750}{\celsius}) closes first, hornblende (about
\qty{530}{\celsius}) next, biotite (about \qty{280}{\celsius}) last, so the ages
decrease in that order as the granite cooled. (b)~$(530 - 280)/(410 - 398) =
250/12 = \qty{21}{\celsius}$ per million years. (c)~Zircon closed above
\qty{750}{\celsius}, so the granite cooled by at least $750 - 530 =
\qty{220}{\celsius}$ in \qty{10}{\Myr}: at least \qty{22}{\celsius} per million
years, a lower bound.
\end{solution}

\begin{solution}{exo:b1:radiometric-methods:5}
(a)~$\bar x = 7.5$, $\bar y = 0.809525$, $S_{xx} = 81.0$, $S_{xy} = 1.13965$;
$m = 1.13965/81.0 = 0.0140698$, $b = 0.809525 - 7.5 \times 0.0140698 = 0.70400$.
(b)~$t = \ln(1.0140698)/\num{1.3972e-11} = \qty{1000.0}{\Ma}$; \omterm{def:b1:radiometric-methods:isochron}{initial ratio}
$0.7040$. (c)~The residuals are $+1.6$, $-2.5$, $+0.7$ and $+0.2$ (in units of
\num{e-4}): $s = \sqrt{\num{9.34e-8}/2} = 0.00022$, $\sigma_m = 0.00022/\sqrt{81}
= \num{2.40e-5}$ and $\sigma_t = \num{2.40e-5}/(\num{1.3972e-11} \times 1.0141) =
\qty{1.7}{\Myr}$: \qty[uncertainty-mode=separate]{1000.0(17)}{\Ma}. (d)~$0.7040$
is close to the \omterm{def:b1:layered-earth:layers}{mantle} value $0.7025$ and well below the crustal $0.710$: a source
mostly in the \omterm{def:b1:layered-earth:layers}{mantle}.
\end{solution}

\begin{solution}{exo:b1:radiometric-methods:6}
(a)~$0.100/39.098 \times \num{1.17e-4} = \qty{2.992e-7}{mol}$ of potassium-40 per
gram. (b)~${}^{40}\mathrm{Ar}^*/{}^{40}\mathrm{K} = \num{1.169e-9}/\num{2.992e-7}
= 0.0039065$; $\lambda/\lambda_e = 5.543/0.581 = 9.5404$; $t = \ln(1 + 9.5404
\times 0.0039065)/\num{5.543e-10} = \ln(1.037269)/\num{5.543e-10} =
\qty{66.01}{\Ma}$. (c)~The radiogenic argon would be $0.995 \times 1.169$
nanomoles, giving \qty{65.69}{\Ma}: the measured age would be \qty{0.32}{\Myr}
too old, ten times the \num{32000} years within which the impact and the
extinction at the end of the Cretaceous are now shown to coincide. Excess argon is a concern mostly in young rocks and whole rocks; a
\omterm{def:b1:rock-forming-minerals:mineral}{mineral} such as sanidine should not take in argon when it crystallises, and an
argon--argon age spectrum would reveal it.
\end{solution}

\begin{solution}{exo:b1:radiometric-methods:7}
(a)~About \num{2000} cal BP: the curve passes through \num{2070} years BP there.
(b)~From about \num{2495} to \num{2699} cal BP, two hundred years: the curve is
nearly flat between \num{2500} and \num{2700} cal BP, so that every calendar age
in that interval has nearly the same \omterm{def:b1:radiometric-methods:radiocarbon}{radiocarbon age}, and a more precise
measurement would not narrow the range. (c)~At \num{2000} cal BP the \omterm{def:b1:radiometric-methods:radiocarbon}{radiocarbon
age} is \num{70} years \emph{older} than the calendar age; at \num{8000} cal BP
\num{822} years younger (\qty{10}{\percent}); at \num{12000} cal BP \num{1694}
years younger (\qty{14}{\percent}). The difference is not proportional to the
age: it even changes sign. It cannot come from the half-life alone, which would
give a constant \qty{2.4}{\percent}.
\end{solution}

\begin{solution}{exo:b1:radiometric-methods:8}
(a)~The ash bed below (\qty{94.50}{\Ma}) gives a maximum age, the lava flow above
(\qty{93.80}{\Ma}) a minimum: the boundary is bracketed between them. The dyke is
also a minimum, but a weaker one, younger than the lava. The detrital zircon is a
maximum only for the sandstone, and weaker than the ash bed, which lies below the
sandstone. (b)~$f = (12 - 4)/(20 - 4) = 0.5$; $t = 94.50 - 0.5 \times 0.70 =
\qty{94.15}{\Ma}$; $\sigma_t = \sqrt{0.25 \times 0.10^2 + 0.25 \times 0.12^2} =
\qty{0.08}{\Myr}$: \qty[uncertainty-mode=separate]{94.15(8)}{\Ma}. (c)~Its
grains came from a rock at least \qty{95.3}{\Ma} old, eroded somewhere in the
source area; it shows that the sandstone is younger than \qty{95.3}{\Ma}, which
the ash bed below already proves more tightly.
\end{solution}

\begin{solution}{exo:b1:radiometric-methods:9}
(a)~Grain~A: $\ln(1.08401)/\num{1.55125e-10} = \qty{520.0}{\Ma}$ and
$\ln(1.6688)/\num{9.8485e-10} = \qty{520.0}{\Ma}$. Grain~B: \qty{368.4}{\Ma} and
\qty{389.9}{\Ma}. (b)~Grain~A is concordant. (c)~$0.05881/0.08401 = 0.7000$ and
$0.4682/0.6688 = 0.7001$. The age is $\ln(1 + R)/\lambda$, not proportional to
$R$: for the small uranium-238 ratio $\ln(1+R) \approx R$ and the age falls by
nearly \qty{30}{\percent} (to \qty{368}{\Ma}); for the larger uranium-235 ratio
the logarithm falls by less, and the age only to \qty{390}{\Ma}. Both are too
young, and they disagree: the grain is discordant.
\end{solution}

\begin{solution}{exo:b1:radiometric-methods:10}
(a)~With $u = \lambda t$, $D^*/P = e^u - 1 = u + u^2/2 + \dots$, so $(D^*/P)/\lambda
= t\,(1 + u/2 + \dots)$: the approximation overestimates the age by a relative
error close to $u/2$. (b)~Rubidium-87 at \qty{310}{\Ma}: $u = 0.00433$, error
\qty{0.22}{\percent} (\qty{0.7}{\Myr}); potassium-40 at \qty{296}{\Ma}: $u =
0.164$, error \qty{8.7}{\percent} (\qty{26}{\Myr}; $u/2$ gives \qty{8.2}{\percent});
uranium-238 at \qty{4.4}{\Ga}: $u = 0.683$, error \qty{43}{\percent} --- the
approximation and even its error estimate ($u/2 = \qty{34}{\percent}$) fail.
(c)~$R = 0.004341$, $\sigma_R = \num{4.34e-6}$; by
\cref{prop:b1:radiometric-methods:precision}, $\sigma_t =
\num{4.34e-6}/(\num{1.3972e-11} \times 1.0043) = \qty{0.31}{\Myr}$.
\end{solution}

\begin{solution}{exo:b1:radiometric-methods:11}
(a)~Let ${}^{40}\mathrm{K}$ be today's potassium-40. At the time $t_2$ before
the present there were ${}^{40}\mathrm{K}\,e^{\lambda t_2}$ atoms, and the argon
made between $t_1$ and $t_2$ was, by \cref{prop:b1:radiometric-methods:kar} applied
to that interval, $(\lambda_e/\lambda)\,{}^{40}\mathrm{K}\,e^{\lambda
t_2}\bigl(e^{\lambda (t_1 - t_2)} - 1\bigr) = (\lambda_e/\lambda)\,
{}^{40}\mathrm{K}\,\bigl(e^{\lambda t_1} - e^{\lambda t_2}\bigr)$. A fraction
$1 - f$ of it remained. Since $t_2$, the closed \omterm{def:b1:rock-forming-minerals:mineral}{mineral} has made
$(\lambda_e/\lambda)\,{}^{40}\mathrm{K}\,(e^{\lambda t_2} - 1)$ more. Adding and
multiplying by $\lambda/(\lambda_e\,{}^{40}\mathrm{K})$ gives the expression.
(b)~The apparent age is $\ln(1 + \text{expression})/\lambda$: \qty{400.0}{\Myr} for
$f = 0$, \qty{342.3}{\Myr} for $0.3$, \qty{282.7}{\Myr} for $0.6$ and
\qty{200.0}{\Myr} for $1$. (c)~Every age between \qty{200}{\Myr} and
\qty{400}{\Myr} can be produced by some partial loss, and a single number cannot
say whether it records a real event at that time or a mixture; an age spectrum,
or other clocks in the same rock, are needed.
\end{solution}

\begin{solution}{exo:b1:radiometric-methods:12}
(a)~The age is proportional to the mean life used: $10306 \times 8223/8033 =
\num{10550}$ years. (b)~$10550 - 10306 = 244$ years, about a seventh of the
\num{1694}. (c)~If the initial activity was $k A_0$, the true age is $t = \tau
\ln(k A_0/A) = \num{10550} + \tau \ln k$; with $t = \num{12000}$, $\ln k =
1450/8223 = 0.176$ and $k = 1.19$: about \qty{19}{\percent} more carbon-14 than
in the standard.
\end{solution}

\begin{solution}{pb:b1:radiometric-methods:1}
\begin{enumerate}
  \item The points rise steadily with $x$; their deviations from a line are at most
        two units in the fourth decimal, against a range of $0.04$ in $y$: they lie on a
        straight line within small scatter.
  \item $\bar x = 4.90333$, $\bar y = 0.729288$, $S_{xx} = 62.1253$, $S_{xy} =
        0.269335$; $m = 0.269335/62.1253 = 0.0043354$.
  \item $b = 0.729288 - 0.0043354 \times 4.90333 = 0.70803$: the
        $^{87}\mathrm{Sr}/^{86}\mathrm{Sr}$ of the granite when it crystallised,
        shared by all its parts.
  \item $t = \ln(1.0043354)/\num{1.3972e-11} = \qty{309.6}{\Ma}$.
  \item The residuals are $+9.1$, $-17.2$, $+6.9$, $-6.4$, $+15.1$ and $-7.5$ (in
        units of \num{e-5}); $s = 0.000137$, $\sigma_m = 0.000137/\sqrt{62.1253} =
        \num{1.74e-5}$, $\sigma_t = \num{1.74e-5}/(\num{1.3972e-11} \times 1.0043)
        = \qty{1.2}{\Myr}$: \qty[uncertainty-mode=separate]{309.6(12)}{\Ma}, and
        $\sigma_b = 0.00010$.
  \item $0.7080$ is closer to the crustal $0.710$ than to the \omterm{def:b1:layered-earth:layers}{mantle}'s $0.7025$.
        For a mixture, $(0.70803 - 0.7025)/(0.710 - 0.7025) = 0.74$: about three
        quarters of the strontium came from melted \omterm{def:b1:layered-earth:layers}{crust}.
  \item $0.0740/39.098 \times \num{1.17e-4} = \qty{2.214e-7}{mol}$ per gram,
        times $\num{6.022e23}$: \num{1.33e17} atoms.
  \item $^{40}\mathrm{Ar}^*/^{40}\mathrm{K} = \num{4.139e-9}/\num{2.2144e-7} =
        0.018691$; $1 + 9.5404 \times 0.018691 = 1.17832$; $t =
        \ln(1.17832)/\num{5.543e-10} = \qty{296.0}{\Ma}$.
  \item $^{40}\mathrm{K} = \qty{3.292e-8}{mol.g^{-1}}$, ratio $0.019239$,
        $\ln(1.18355)/\num{5.543e-10} = \qty{304.0}{\Ma}$.
  \item $\ln(1.018691)/\num{5.543e-10} = \qty{33.4}{\Ma}$, nine times too young,
        because only about one decay in ten gives argon. Calcium-40 is the common
        isotope of calcium, abundant in rocks: the little made by decay cannot be
        told from the calcium the \omterm{def:b1:rock-forming-minerals:mineral}{minerals} contained from the start.
  \item Excess argon shifts ages by amounts of the order of a million years, which
        matters for young rocks; here it would be well under \qty{1}{\percent} of
        the age. And \omterm{def:b1:rock-forming-minerals:groups}{mica} and \omterm{def:b1:rock-forming-minerals:groups}{amphibole} crystals take in almost no argon when they
        form.
  \item Either the granite cooled slowly and the \omterm{def:b1:rock-forming-minerals:groups}{micas} and \omterm{def:b1:rock-forming-minerals:groups}{amphibole} closed after
        crystallisation (\omterm{def:b1:radiometric-methods:closure}{cooling ages}), or a later heating drove out part of their
        argon. A flat plateau in an argon--argon age spectrum supports a \omterm{def:b1:radiometric-methods:closure}{cooling
        age}; a spectrum without a plateau warns that the \omterm{def:b1:rock-forming-minerals:mineral}{minerals} lost argon. The
        order of the ages, hornblende before biotite as their \omterm{def:b1:radiometric-methods:closure}{closure
        temperatures} predict, also points to cooling.
  \item $\ln(1.049361)/\num{1.55125e-10} = \qty{310.6}{\Ma}$ and
        $\ln(1.35784)/\num{9.8485e-10} = \qty{310.6}{\Ma}$: concordant.
  \item Zircon \qty{310.6}{\Ma}, closure above \qty{750}{\celsius}:
        crystallisation; the whole-rock \omterm{def:b1:radiometric-methods:isochron}{isochron},
        \qty[uncertainty-mode=separate]{309.6(12)}{\Ma}, also dates
        it, and agrees within its uncertainty; hornblende \qty{304.0}{\Ma}: cooling
        through about \qty{530}{\celsius}; biotite \qty{296.0}{\Ma}: cooling
        through about \qty{280}{\celsius}. The graph falls from above
        \qty{750}{\celsius} at \qty{310.6}{\Ma} to \qty{280}{\celsius} at
        \qty{296.0}{\Ma}.
  \item $(530 - 280)/(304.02 - 296.03) = 250/7.99 = \qty{31}{\celsius}$ per
        million years.
  \item $(750 - 530)/(310.60 - 304.02) = 220/6.58 = \qty{33}{\celsius}$ per
        million years at least, since the zircon closed above
        \qty{750}{\celsius}: the cooling was at least as fast early on.
  \item The difference of temperatures, $250 \pm \sqrt{40^2 + 40^2} = 250 \pm
        57$, gives \qtyrange{24}{38}{\celsius} per million years.
  \item $(550 - 350)/7.99 = \qty{25}{\celsius}$ per million years, inside the range
        of question~17: the uncertainty comes from the \omterm{def:b1:radiometric-methods:closure}{closure temperatures}, not
        from the ages.
  \item $(280 - 10)/30 = \qty{9.0}{km}$.
  \item $\ln(1.039108)/\num{1.55125e-10} = \qty{247.30}{\Ma}$, and the same from
        the uranium-235 ratio: concordant.
  \item From the granite, eroded at the \omterm{def:b1:cross-sections:unconformity}{unconformity}: their age is its
        crystallisation age. They show only that the sandstone is younger than
        \qty{310}{\Ma}: a maximum age, much weaker than the ash bed.
  \item $(58 - 18)/(247.30 - 245.30) = \qty{20.0}{m}$ per million years; the base
        is $18/20.0 = \qty{0.90}{\Myr}$ older than the lower ash:
        \qty{248.20}{\Ma}.
  \item From \qty{309.6}{\Ma} (\omterm{def:b1:radiometric-methods:isochron}{isochron}) or \qty{310.6}{\Ma} (zircon) to
        \qty{248.2}{\Ma}: about \qty{62}{\Myr}.
  \item \qty{9.0}{km} in $296.0 - 248.2 = \qty{47.8}{\Myr}$: \qty{0.19}{km} per
        million years, \qty{0.19}{mm} per year. This assumes that the dated
        biotite came from just below the \omterm{def:b1:cross-sections:unconformity}{unconformity}, so that all the rock above
        it in the granite and the schist was removed by the time the sand arrived,
        and that the \omterm{def:b1:heat-budget:geotherm}{geotherm} stayed constant.
  \item The granite cooled at about \qty{31}{\celsius} per million years
        between the closures of hornblende and biotite.
\end{enumerate}
\end{solution}
